\chapter{Experiments}\label{ch:experiments}

This chapter reports synthetic experiments that test the algorithm of
Chapter~\ref{ch:algorithms} under controlled conditions. The data are generated
from the model of Chapter~\ref{ch:model} so that the true parameters are known.

\section{Setup}\label{sec:setup}

We generate samples of size $n\in\{200,1000,5000\}$ in dimension $d=2$ with
$k=3$ components. Each configuration is repeated with twenty random seeds, and
the initialisation strategies of Table~\ref{tab:init} are compared.

\lipsum[35-36]

\section{Recovery accuracy}\label{sec:recovery}

Table~\ref{tab:recovery} reports the mean absolute error in the estimated
mixing weights, averaged over seeds. The errors decrease with $n$, as expected
from the usual asymptotic theory.

\begin{table}[ht]
  \centering
  \begin{tabular}{rccc}
    \toprule
    $n$ & Random & $k$-means & Hierarchical \\
    \midrule
    200  & 0.081 & 0.047 & 0.052 \\
    1000 & 0.034 & 0.021 & 0.023 \\
    5000 & 0.012 & 0.009 & 0.009 \\
    \bottomrule
  \end{tabular}
  \caption{Mean absolute error in mixing weights across initialisation strategies.}
  \label{tab:recovery}
\end{table}

\lipsum[37-38]

\section{Iteration counts}\label{sec:iterations}

The number of EM iterations needed to meet the stopping rule of
Equation~\eqref{eq:stop} is reported in Section~\ref{sec:discussion-speed}
of the discussion. Initialisation by $k$-means typically halves the count
compared with random starts.

\lipsum[39-40]

\section{Further remarks}\label{sec:further-experiments}

The remarks in this section restate the main point of the chapter in more
general terms, so that the argument can be read without the surrounding
derivations.

\lipsum[111-122]
